%% kapa: exam=1 round=27 year=2016 period=2 form=A subject=appraisal_law subject_ko=감정평가관계법규 seq=35 number=35
%% source: Q-Net 감정평가사 기출문제 (cst00309) artlSeq=5169402 file=extracted/1008/5169402/제27회 감정평가사 1차 문제/제27회 감정평가사 2교시문제(A형).pdf
%% conversion: kordoc 4.24.51 -> tools/gen1.py (auto); review=proofread-sonnet; flags=none
\begin{question}{35}
부동산등기법령상 등기관의 처분에 대한 이의에 관한 설명으로 옳은 것은?
\begin{choices}
  \item 이의의 신청은「민사소송법」이 정하는 바에 따라 관할 지방법원에 이의신청서를 제출하는 방법으로 한다.
  \item 새로운 사실이나 새로운 증거방법을 근거로 이의신청을 할 수 있다.
  \item 등기관은 이의가 이유 있다고 인정하더라도 그에 해당하는 처분을 해서는 아니되고 관할 지방법원에 보내 그 결정에 따라야 한다.
  \item 이의에는 집행정지의 효력이 없다.
  \item 이의에 대한 관할 지방법원의 결정에 대해서는 불복할 수 없다.
\end{choices}
\end{question}
\answer{4}
